\chapter{Introduction}
\label{chap:introduction}

Higher education increasingly relies on digital infrastructure — but most universities still run practical laboratory examinations the old way \cite{pmc_cbt}. Schedules are shared on WhatsApp groups. Attendance is marked on paper. Students email their code files at the end. Supervisors watch thirty screens with two eyes. This is not a scalable process.

Computer science and software engineering programs depend on practical exams more than most. These exams test actual coding ability, system design judgment, and real-time problem solving under time pressure. They cannot be reduced to a multiple-choice quiz. They need a working environment — a machine, an IDE, a compiler — and a supervisor who knows what is happening on every screen simultaneously.

The proposed AI-Driven Lab Exam Monitoring and Management System targets this exact gap. It integrates scheduling, identity verification, workstation lockdown, live telemetry monitoring, auto-saved submissions, AI-assisted grading, and analytical PDF reporting into one auditable platform \cite{pmc_cbt}.

The problem with the current state is not that supervisors lack effort. It is that the tools are fragmented. Scheduling happens in one place, attendance in another, file collection in a third. None of these tools talk to each other. Evidence is scattered. When a student disputes a result or a suspicious incident occurs during the exam, there is no single source of truth to consult.

A unified system changes this. Administrators configure access centrally. Teachers build exams — with allowed application lists, test cases, and time limits — before the session begins. Students log in through a hardware-bound client that locks the desktop to the exam environment. Invigilators monitor a live grid dashboard showing each student's focus timeline and process activity. After the exam ends, the system generates a structured report \cite{integ1}. The audit chain is complete.

SmartExam is that unified system.
